% Appendix H -- REWRITTEN 2026-10-06. Was "The superseded rider sweep": the
% four-point sweep set aside in 2026-09-26, kept so a reader who had met its 78%
% overlap figure could see what replaced it.
%
% That framing no longer fits, because the thing that replaced it has itself been
% replaced. The four-point sweep was not merely imprecise, it was fitted ACROSS the
% knee, and so was its successor. This appendix now carries the characterisation
% done below the knee, the coefficient table with provenance, and the full ledger
% of optimisations attempted -- including the ones that failed, which is the part
% a reader extending this work most needs.
%
% The long-form version, with every arm and every retraction, is
% paper/OVERHEAD_CHARACTERISATION.md in the repository; it is not part of the
% paper.

\section{The overhead law, measured}
\label{app:ridercost}

\subsection{Why the earlier fits were wrong, and in which direction}

Three successive readings of the per-rider cost appear in this project's history,
and all three errors have one cause: \textbf{points above the knee were included in
the fit}. Above $r^{*}$ (Eq.~\ref{eq:knee}) the step time \emph{is} the transfer, so
it grows with $r$ at the link rate rather than at the overhead rate. A line fitted
through points on both sides of the boundary therefore tilts --- it borrows slope
from the transfer-bound regime and pays for it with a depressed intercept.

That single mechanism explains every discrepancy we chased:

\begin{itemize}\itemsep1pt \parskip0pt
\item the per-rider cost came out $2.2$--$2.9$\,ms, against $0.52$ measured below
  the knee: \textbf{over} by four to five times;
\item the fixed cost came out $9.6$--$12.3$\,ms, against $3.4$: \textbf{under},
  once the optimisations below are also accounted for;
\item an apparent \emph{negative} fixed cost at one configuration, which we
  reported as evidence that the law was not affine;
\item an apparent dependence of the per-rider cost on context,
  $c = 0.205 + 1.84\times10^{-3}S$, which vanishes entirely below the knee;
\item an apparent $62\%$, later $88\%$, transfer-overlap efficiency --- a quantity
  that does not exist, since below the knee the overlap is complete.
\end{itemize}

The lesson is procedural rather than numerical. The knee must be computed
\emph{per point} from the final sequence length, not the prompt length, and points
failing $\ell/T_0 \le 1/3$ must be discarded rather than down-weighted. Doing that
is what makes the two coefficients separately identifiable: across seven arms the
per-rider cost repeats to $1.5\%$ and the fixed cost to $10\%$ between runs of one
configuration.

\subsection{The coefficients}

\begin{center}\small
\begin{tabular}{llr}
\toprule
Quantity & Value & Provenance \\
\midrule
$c$ intercept       & $0.177$\,ms            & $L = 16, 32, 64$ at ctx 128 \\
$c$ per layer       & $0.0054$\,ms           & same \\
$F$ per layer       & $0.053$\,ms            & after T1, T2, T7a, T8 \\
mirror write        & $1.2$--$1.5$\,ms       & ablation at $L = 64$ \\
second attention    & $1.9$\,ms              & ablation at $L = 64$ \\
D2H copy floor      & $0.834$\,ms            & off-engine, $2L$ copies \\
\bottomrule
\end{tabular}
\end{center}
% src: paper/data/overhead/15_clean_cf/CF_L16.json,
% paper/data/overhead/15_clean_cf/CF_L32.json,
% paper/data/overhead/15_clean_cf/CF_c128_a.json,
% paper/data/overhead/20_f_split_post_t8/FS_nod2h.json,
% paper/data/overhead/20_f_split_post_t8/FS_both.json

Two cautions. The fixed cost's run-to-run drift on identical configuration was
$0.04$, $0.43$, $0.04$ and $0.51$\,ms across four jobs --- \textbf{bimodal}, so a
tight bracket in one job licenses nothing about another; we quote $\pm 0.5$\,ms and
re-measure the baseline inside every job. And the two components of $F$ are each a
single arm, so they are $\pm 0.5$ individually while their sum is pinned by the
tighter three-arm baseline.

\subsection{The ledger: nine attempted, four landed}

\begin{center}\small
\begin{tabular}{llp{4.6cm}}
\toprule
 & Change & Outcome \\
\midrule
T1 & $O(r)$ extent allocation   & \textbf{landed}: $c$ $5.17 \to 1.26$\,ms \\
T2 & deferred batched mirror    & \textbf{landed}: $F$ $13.4 \to 5.8$\,ms \\
T7a & vectorised mirror scatter & \textbf{landed}: $-0.32$\,ms \\
T8 & stop a disabled feature flushing per layer & \textbf{landed}: $-2.0$\,ms \\
T3 & staged D2D then D2H        & rejected: $6\times$ worse, 4.6 vs 22.3\,GB/s \\
T4 & batch copies across layers & rejected unbuilt: $18\%$ worse off-engine \\
T5 & whole-extent runs          & null: launches $\div 8$, no resolvable change \\
T6 & order the mirror scatter   & rejected: costs the entire overlap \\
T7b & preallocated rider buffers & null \\
\bottomrule
\end{tabular}
\end{center}
% src: paper/data/overhead/08_opt/T1_c2048_bs16.json,
% paper/data/overhead/09_t2/T2_async_a.json,
% paper/data/overhead/10_t3/T3_on_a.json,
% paper/data/overhead/14_t5/T5_on.json,
% paper/data/overhead/13_t6/T6_on.json,
% paper/data/overhead/19_t8/T8_on_a.json

\textbf{Four of the nine had a mechanism that demonstrably fired and a benefit that
did not follow}, and that ratio is the most transferable result in this appendix. T5
cut the launch count eightfold and moved nothing, which is what refuted the
launch-bound account of $F$. T3's staged copy achieved exactly the bandwidth
improvement predicted for it in isolation and was six times slower in the engine. We
therefore price every remaining candidate off-engine or by ablation before building
it, and treat a confirmed mechanism as no evidence at all about a cost.

The largest single win was not an optimisation. A feature that is \textbf{off by
default} --- a swap-semantics emulation --- was flushing the batched mirror write on
every layer, $64$ times per step, because its call site tested that the method
\emph{existed} rather than that the feature was enabled. It had silently defeated
T2's batching from the day T2 landed. It was found only because T7a, a change that
should have helped, did not: the diagnostic counter added to explain T7a's silence
reported a state that contradicted the row counts.

\subsection{What remains, and its measured ceiling}

$56\%$ of the fixed cost is now the second attention call that the two-source design
requires (\S\ref{sec:design:exec}). Removing it means fusing the two calls into one
--- a single varlen invocation over a concatenated block table, which needs the rider
cohort's table rebuilt per layer on the hot path. Given the ledger above we declined
to project the benefit and measured its ceiling instead, by ablating the call
entirely: a fused implementation cannot beat not doing the work at all. The remaining
mirror write is bounded below by $0.834$\,ms, the cost of its $2L$ copies measured
off-engine, so at most $0.6$\,ms of it is recoverable.
